\section{Introduction}

% Paragraph 1: Introduce DeFi networks, credit exposures, risks, and measurement gaps
Decentralised finance (DeFi) has become a significant and rapidly expanding segment of the crypto‑asset ecosystem, locking billions of dollars in value and promising greater financial inclusion, innovation and efficiency \cite{Adamyk2025}.  DeFi networks comprise a rapidly expanding array of protocols—lending platforms, decentralised exchanges, derivatives, asset‑management products and stablecoins—whose interactions are governed by smart contracts and token flows.  Unlike traditional finance, where credit exposure is created by explicit bilateral contracts, DeFi exposure arises when one protocol holds or locks tokens issued by another; this token‑mediated dependence creates a complex network of interdependencies \cite{DeXposure2025}.  A protocol’s tokens may be adopted as collateral, liquidity or reserves across multiple chains, so the failure of one protocol can propagate contagion to all that hold its tokens \cite{DeXposure2025}.  Recent research using the DeXposure dataset shows that inter‑protocol credit exposure has grown rapidly, is concentrated in a small set of core protocols and exhibits declining network density, highlighting a trend toward systemic concentration \cite{DeXposure2025}.  Yet most empirical studies rely on static snapshots or small sets of lending protocols, so there is a critical need for dynamic, ecosystem‑wide datasets and models to monitor how exposure concentrations shift during market events and how new exposure pathways emerge \cite{DeXposure2025}.

The Financial Stability Board warns that by attempting to replicate traditional financial functions, DeFi inherits and may amplify vulnerabilities such as operational fragilities, liquidity and maturity mismatches, leverage and interconnectedness \cite{FSB2023}.  These vulnerabilities can play out differently because automated liquidation mechanisms and reliance on oracles may trigger rapid deleveraging and fire sales \cite{FSB2023}.  The European Systemic Risk Board notes that, by mid‑2025, links between the crypto sector and the broader financial system had intensified, raising financial‑stability concerns, particularly through spillover channels associated with stablecoins \cite{ESRB2025}.  Academic analyses further underscore that DeFi’s benefits—such as financial inclusion and operational efficiency—are accompanied by smart‑contract vulnerabilities, price volatility, market manipulation and liquidity crises; users have suffered significant losses from flash‑loan attacks and coding errors \cite{Adamyk2025}.  Because assets are locked across multiple protocols, a single failure can cascade through the ecosystem, making early detection of vulnerabilities and robust, real‑time risk management essential \cite{Adamyk2025}.  Despite these mounting risks and the growth of inter‑protocol credit exposures, the systemic implications of DeFi networks remain poorly understood due to data limitations and a lack of standardised risk metrics.  A deeper understanding of these exposures is vital for assessing stability, designing appropriate regulation and developing tools to manage and mitigate contagion in decentralised financial markets.


% Paragraph 2: Introduce the DeXposure dataset and DeXposure‑FM model to address the issue
To address these gaps, we build on the **DeXposure** dataset—a large‑scale, publicly released dataset that reconstructs value‑linked credit exposures across more than 4,300+ protocols, 602 blockchains and 24,300+ tokens between 2020 and 2025  \cite{wu2025dexposure}.  DeXposure defines inter‑protocol exposures through changes in total value locked (TVL) and uses DefiLlama metadata to infer token flows, providing the first comprehensive measurement of inter‑protocol credit dependencies.  On top of this dataset we train **DeXposure‑FM**, a multimodal foundation model designed to measure and forecast credit exposures in global DeFi networks.  The model combines time‑series, graph and tabular encoders to jointly represent protocol‑level TVL and flows, dynamic network topology and sector attributes.  Training employs **Sharpness‑Aware Minimization** (SAM)—an optimization procedure that seeks parameters in flat regions of the loss landscape—to improve generalization  \cite{oikonomou2025sam}.  By leveraging foundation‑model principles, including modality‑aware architectures and pre‑training strategies  \cite{eremeev2025g2tfm,liang2024fmts}, DeXposure‑FM can incorporate heterogeneous modalities and long temporal horizons while preserving coherence across modalities.

% Paragraph 3: Describe the machine‑learning benchmarks used to evaluate the model
We evaluate DeXposure‑FM on a suite of machine‑learning tasks that reflect both policy relevance and methodological challenge.  First, the model forecasts edge‑level exposures and network‑level statistics such as concentration, density and sector connectivity several steps ahead.  Second, we perform stress‑testing scenario analysis by generating shock‑consistent paths of losses and contagion during events like stablecoin de‑pegs or exchange failures.  Third, we assess the model’s ability to impute and denoise missing or noisy exposure and TVL histories, using TVL‑weighted error metrics and downstream performance on forecasting and link‑prediction tasks.  These benchmarks extend earlier work on vector autoregression, low‑rank factor models and temporal graph neural networks and demonstrate the benefits of a multimodal foundation model.

% Paragraph 4: Explain financial‑economics insights generated by the model
Beyond machine‑learning benchmarks, DeXposure‑FM delivers financial economics insights with direct macroprudential relevance.  The model yields dynamic measures of protocol‑level systemic importance, sector‑to‑sector spillover indices and early‑warning indicators based on shifts in network concentration and dependence on fragile collateral.  These statistics respond to policy calls for **AI‑driven economic measurement**   \cite{nber2025ai} and illustrate how multimodal foundation models trained on DeFi data can support macroprudential monitoring, stress testing and counterfactual policy evaluation while highlighting the conditions under which the resulting measures are robust or warrant caution.  In doing so, DeXposure‑FM contributes to the emerging literature on systemic risk in DeFi  \cite{naifar2025mapping,zbandut2025institutionalizing} and demonstrates the potential of foundation models for economic measurement and financial regulation.

The model is available at ; the training and deployment code is released at \texttt{https://github.com/EVIEHub/DeXposure-FM}.